\chapter{附录}
\section{特征值问题}
以下Robin条件中的$h,h_1,h_2$均取正值；若允许负值，还须另外讨论零和负特征值。
\begin{enumerate}
	\item $\begin{cases}
		X'' + \lambda X = 0\\
		X(0) = X(L) = 0
	\end{cases}$
	\begin{equation}
		\label{eq:mmp-interval-dirichlet-spectrum}
		\lambda_n = \Bigp{\frac{n\pi}L}^2,\quad X_n = \sin\frac{n\pi x}L,\quad n = 1,2,\ldots.
	\end{equation}
	\item $\begin{cases}
		X'' + \lambda X = 0\\
		X'(0) = X'(L) = 0
	\end{cases}$
	\begin{equation}
		\label{eq:mmp-interval-neumann-spectrum}
		\lambda_n = \Bigp{\frac{n\pi}L}^2,\quad X_n = \cos\frac{n\pi x}L,\quad n = {\color{red}0},1,\ldots.
	\end{equation}
	\item $\begin{cases}
		X'' + \lambda X = 0\\
		X'(0) = X(L) = 0
	\end{cases}$
	\begin{equation}
		\label{eq:mmp-interval-neumann-dirichlet-spectrum}
		\lambda_n = \Bigb{\Bigp{n + \frac12}\frac{\pi}L}^2,\quad X_n = \cos\Bigp{n + \frac12}\frac{\pi x}L,\quad n = 0,1,\ldots.
	\end{equation}
	\item $\begin{cases}
		X'' + \lambda X = 0\\
		X(0) = X'(L) = 0
	\end{cases}$
	\begin{equation}
		\label{eq:mmp-interval-dirichlet-neumann-spectrum}
		\lambda_n = \Bigb{\Bigp{n + \frac12}\frac{\pi}L}^2,\quad X_n = \sin\Bigp{n + \frac12}\frac{\pi x}L,\quad n = 0,1,\ldots.
	\end{equation}
	\item $\begin{cases}
		X'' + \lambda X = 0\\
		X(0) = X'(L) + hX(L) = 0
	\end{cases} \implies \gamma = -hL\tan\gamma > 0,$
	\begin{equation}
		\label{eq:mmp-interval-dirichlet-robin-spectrum}
		\lambda_n = \Bigp{\frac{\gamma_n}L}^2,\quad X_n = \sin\frac{\gamma_n x}L,\quad n = 1,2,\ldots.
	\end{equation}
	\item $\begin{cases}
		X'' + \lambda X = 0\\
		X'(0) = X'(L) + hX(L) = 0
	\end{cases} \implies \gamma = hL\cot\gamma > 0,$
	\begin{equation}
		\label{eq:mmp-interval-neumann-robin-spectrum}
		\lambda_n = \Bigp{\frac{\gamma_n}L}^2,\quad X_n = \cos\frac{\gamma_n x}L,\quad n = 1,2,\ldots.
	\end{equation}
	\item $\begin{cases}
		X'' + \lambda X = 0\\
		X(0) - h_1X'(0) = 0\\
		X(L) + h_2X'(L) = 0
	\end{cases} \implies (1 - h_1h_2\gamma^2)\sin(L\gamma) + (h_1 + h_2)\gamma\cos(L\gamma) = 0,\quad\gamma > 0,$
	\begin{equation}
		\label{eq:mmp-interval-robin-spectrum}
		\lambda_n = \gamma_n^2,\quad X_n = \sin\gamma_n x + h_1\gamma_n\cos\gamma_nx,\quad n = 1,2,\ldots.
	\end{equation}
	\item $\begin{cases}
		\varTheta'' + \lambda\varTheta = 0\\
		\varTheta(\theta + 2\pi) = \varTheta(\theta)
	\end{cases}$
	\begin{equation}
		\label{eq:mmp-angular-periodic-spectrum}
		\lambda_n = n^2,\quad\varTheta_n(\theta) = A\cos n\theta + B\sin n\theta,\quad n = 0,1,\ldots.
	\end{equation}
\end{enumerate}

\section{Fourier和Laplace}

\paragraph{Fourier级数}
定义实函数$f,g$在区间$[-\pi,\pi]$上的内积
\begin{equation}
	\label{eq:mmp-fourier-inner-product}
	\inp fg := \int_{-\pi}^\pi f(x)g(x)\d x.
\end{equation}
对正整数$m,n$，注意到
\begin{subequations}
	\label{eqn:intsincos}
	\begin{align}
		\int_{-\pi}^\pi\sin nx\d x &= \int_{-\pi}^\pi\cos nx\d x = 0.\\
		\int_{-\pi}^\pi\cos nx\cos mx\d x & = \frac12\int_{-\pi}^\pi\cos(n - m)x + \cos(n + m)x\d x = \pi\delta_{nm};\\
		\int_{-\pi}^\pi\sin nx\cos mx\d x & = \frac12\int_{-\pi}^\pi\sin(n + m)x + \sin(n - m)x\d x = 0; \\
		\int_{-\pi}^\pi\sin nx\sin mx\d x & = \frac12\int_{-\pi}^\pi\cos(n - m)x - \cos(n + m)x\d x = \pi\delta_{nm} .
	\end{align}
\end{subequations}
因此
\[\{1,\sin x,\cos x,\sin2x,\ldots\}\]
构成$L^2[-\pi,\pi]$的一组完备正交基，归一化之
\begin{equation}
	\label{eq:mmp-trigonometric-orthonormal-basis}
	\biggB{\frac1{\sqrt{2\pi}},\frac{\sin x}{\sqrt\pi},\frac{\cos x}{\sqrt\pi},\frac{\sin2x}{\sqrt\pi},\ldots}.
\end{equation}

用这一组基可将周期为$2\pi$的平方可积函数$f(x)$按平方平均意义展开
\begin{equation}
	\label{eq:mmp-trigonometric-fourier-series}
	f(x) = \frac{a_0}2 + \sum_{n=1}^\infty(a_n\cos nx + b_n\sin nx),
\end{equation}
其中
\begin{subequations}
	\label{eq:mmp-trigonometric-fourier-coefficients}
	\begin{align}
		a_n & = \frac1\pi\inp f{\cos nx} = \frac1\pi\int_{-\pi}^\pi f(x)\cos nx\d x, \\
		b_n & = \frac1\pi\inp f{\sin nx} = \frac1\pi\int_{-\pi}^\pi f(x)\sin nx\d x.
	\end{align}
\end{subequations}

\paragraph{Fourier变换}
利用$\e^{\i x} = \cos x + \i\sin x$，可得到Fourier级数的复数形式：
\begin{subequations}
	\label{eq:mmp-complex-fourier-series}
	\begin{align}
		f(x) &= \sum_{n=-\infty}^{+\infty}c_n\e^{\i nx},\\
		c_n &= \frac1{2\pi}\int_{-\pi}^\pi f(x)\e^{-\i nx}\d x.
	\end{align}
\end{subequations}
对于其他周期$T = \lambda,\;k = 2\pi/\lambda$
\begin{subequations}
	\label{eq:mmp-periodic-fourier-series}
	\begin{align}
		f(x) &= \sum_{n=-\infty}^{+\infty}c_n\e^{\i nkx},\\
		c_n &= \frac1\lambda\int_{-\lambda/2}^{\lambda/2} f(x)\e^{-\i nkx}\d x.
	\end{align}
\end{subequations}

有限区间上的平方可积函数仍可作Fourier级数展开，其区间外延拓是周期的；若要直接描述整条实轴上的非周期函数，则引入Fourier变换。
然而，若认为其周期$\lambda = 2N$充分大，令$\xi = nk$，
\[f(x) = \sum_\xi\biggb{\frac{k}{2\pi}\int_{-N}^Nf(x)\e^{-\i\xi x}\d x}\e^{\i\xi x}.\]
再让$N \to \infty,\;k \to 0$，定义
\[\hat f(\xi) := \intR f(x)\e^{-\i\xi x}\d x,\]
以及
\[f(x) = \frac1{2\pi}\intR\hat f(\xi)\e^{\i\xi x}\d\xi,\]
这就是Fourier变换。

\subsection{Fourier系数}

下面简要介绍不同函数的Fourier系数如何计算。各不定积分中的 $C$ 表示任意积分常数。

\paragraph{多项式函数}
只需要利用分部积分将$x^k$不断微分降幂，
\begin{subequations}
	\label{eq:mmp-polynomial-trigonometric-primitives}
	\begin{align}
		\int x\cos ax\d x &= \frac{x\sin ax}a + \frac{\cos ax}{a^2} + C.\\
		\int x\sin ax\d x &= -\frac{x\cos ax}a + \frac{\sin ax}{a^2} + C.\\
		\int x^2\cos ax\d x &= \frac{x^2\sin ax}a + \frac{2x\cos ax}{a^2} - \frac{2\sin ax}{a^3} + C.\\
		\int x^2\sin ax\d x &= -\frac{x^2\cos ax}a + \frac{2x\sin ax}{a^2} + \frac{2\cos ax}{a^3} + C.\\
		\int x^3\cos ax\d x &= \frac{x^3\sin ax}a + \frac{3x^2\cos ax}{a^2} - \frac{6x\sin ax}{a^3} - \frac{6\cos ax}{a^4} + C.\\
		\int x^3\sin ax\d x &= -\frac{x^3\cos ax}a + \frac{3x^2\sin ax}{a^2} + \frac{6x\cos ax}{a^3} - \frac{6\sin ax}{a^4} + C.
	\end{align}
\end{subequations}
\begin{remark}可以证明，按照降幂顺序，系数的规律是：
	\begin{itemize}
		\item 正弦系数的符号$-,+,+,-$，余弦系数的符号$+,+,-,-$；
		\item 共有$n + 1$项，系数分别为
		\[
			\frac1a,\enspace\frac n{a^2},\enspace\frac{n(n - 1)}{a^3},\enspace\ldots,\enspace\frac{n!}{a^{n+1}}.
		\]
		\item 正弦函数$\cos$打头，余弦$\sin$打头，然后正余弦交替。
	\end{itemize}
\end{remark}

\paragraph{三角函数}
积分结果已经在\eqref{eqn:intsincos}中给出。

\paragraph{指数函数}
对于指数函数的Fourier系数，只需记住
\begin{subequations}
	\label{eq:mmp-exponential-trigonometric-primitives}
	\begin{align}
		\mathcal I := \int\e^{ax}\cos bx\d x = \e^{ax} \cdot \frac{a\cos bx + b\sin bx}{a^2 + b^2} + C;\\
		\mathcal J := \int\e^{ax}\sin bx\d x = \e^{ax} \cdot \frac{a\sin bx - b\cos bx}{a^2 + b^2} + C.
	\end{align}
\end{subequations}
\begin{proof}
	分部积分
	\begin{align*}
		\mathcal I &= \frac1b\int\e^{ax}\d\sin bx = \frac1b\e^{ax}\sin bx - \frac ab\mathcal J;\\
		\mathcal J &= -\frac1b\int\e^{ax}\d\cos bx = -\frac1b\e^{ax}\cos bx + \frac ab\mathcal I.
	\end{align*}
	即可解得上式。
\end{proof}

\paragraph{双曲函数}
有
\begin{subequations}
	\label{eq:mmp-hyperbolic-trigonometric-primitives}
	\begin{align}
		\int\cosh ax\cos bx\d x = \frac{a\sinh ax\cos bx + b\cosh ax\sin bx}{a^2 + b^2} + C,\\
		\int\cosh ax\sin bx\d x = \frac{a\sinh ax\sin bx - b\cosh ax\cos bx}{a^2 + b^2} + C,\\
		\int\sinh ax\cos bx\d x = \frac{a\cosh ax\cos bx + b\sinh ax\sin bx}{a^2 + b^2} + C,\\
		\int\sinh ax\sin bx\d x = \frac{a\cosh ax\sin bx - b\sinh ax\cos bx}{a^2 + b^2} + C.
	\end{align}
\end{subequations}

\paragraph{对数函数}
对数函数乘正弦函数分部积分后会出现$\Si x$。

\paragraph{反三角函数}
反三角函数更是没有积分结果。

% 因此靠上面的公式可以解决绝大部分Fourier系数的习题。
\subsection{Fourier变换}

\paragraph{Fourier变换的性质}
\label{The Property of Fourier Transformation}
Fourier变换的定义见第~\pageref{dfn:Integral Transformation}~页\dfnref{dfn:Integral Transformation}。

\begin{enumerate}
	\item 线性
		\begin{equation}
			\label{eq:mmp-fourier-linearity}
			\FT[\alpha f + \beta g] = \alpha\FT[f] + \beta\FT[g];
		\end{equation}
	\item 微分
		\begin{equation}
			\label{eq:mmp-fourier-derivative}
			\FT[f'] = \i\xi\hat f.
		\end{equation}
	\item 积分
		\begin{equation}
			\label{eq:mmp-fourier-primitive}
			\FT\biggb{\int_{-\infty}^xf(\eta)\d\eta} = \frac1{\i\xi}\hat f.
		\end{equation}
		其中积分式的Fourier变换存在，且$\hat f(0) = 0$．
	\item $\hat f(\xi)$微分
		\begin{equation}
			\label{eq:mmp-fourier-frequency-derivative}
			\hat f'(\xi) = \FT[-\i xf](\xi).
		\end{equation}
	\item $\hat f(\xi)$积分
			\begin{gather}
				\int_{-\infty}^\xi\hat f(\zeta)\d\zeta = \FT\Bigb{\frac{f(x)}{-\i x}}(\xi) + C_1,\\
				\int_\xi^{+\infty}\hat f(\zeta)\d\zeta = \FT\Bigb{\frac{f(x)}{\i x}}(\xi) + C_2.
			\end{gather}
	\item 位移性质
			\begin{subequations}
				\label{eq:mmp-fourier-shift}
				\begin{gather}
					\FT[f(x - a)](\xi) = \hat f(\xi)\e^{-\i a\xi},\\
					\hat f(\xi - \alpha) = \FT[f(x)\e^{\i\alpha x}](\xi).
				\end{gather}
			\end{subequations}
	\item 相似
			\begin{equation}
				\label{eq:mmp-fourier-dilation}
				\FT[f(kx)](\xi) = \frac1{\abs k}\hat f\Bigp{\frac\xi{k}}.
			\end{equation}
	\item 卷积\footnote{$f,g$的卷积\[f \ast g(x) = \intR f(\eta)g(x - \eta)\d\eta.\]且有$f \ast g = g \ast f$．}
			\begin{equation}
				\label{eq:mmp-fourier-convolution}
				\FT[f \ast g] = \hat f\hat g.
			\end{equation}
	\item 象函数卷积
			\begin{equation}
				\label{eq:mmp-fourier-product}
				\hat f \ast \hat g = 2\pi\FT[fg]
			\end{equation}
	\item 反射
			\begin{equation}
				\label{eq:mmp-fourier-reflection}
				\FT[\hat f(\xi)](x) = 2\pi f(-x).
			\end{equation}
			因此Fourier逆变换
			\[\FT\inv[g](x) = \frac1{2\pi}\intR g(\xi)\e^{\i\xi x}\d\xi.\]
			在$f$连续点处$\FT\inv[\hat f(\xi)](x) = f(x)$．
\end{enumerate}
高维下Fourier变换的性质是相似的．
\paragraph{Fourier正余弦变换的性质}\label{The Property of Sine and Cosine Transformation}\footnote{定义见第~\pageref{dfn:Fourier Sine and Cosine Transformation}~页~\ref{dfn:Fourier Sine and Cosine Transformation}。}
\begin{enumerate}
	\item 微分
	\begin{subequations}
		\label{eq:mmp-half-line-derivatives}
		\begin{align}
			\FT_\sine[f'](\xi) &= f(x)\sin\xi x|_{x=0}^{+\infty} - \xi\FT_\cosine[f](\xi);\\
			\FT_\cosine[f'](\xi) &= f(x)\cos\xi x|_{x=0}^{+\infty} + \xi\FT_\sine[f](\xi);\\
			\FT_\sine[f''](\xi) &= \Bigp{f'(x)\sin\xi x - \xi f(x)\cos\xi x}_{x=0}^{+\infty} - \xi^2\FT_\sine[f](\xi);\\
			\FT_\cosine[f''](\xi) &= \Bigp{f'(x)\cos\xi x + \xi f(x)\sin\xi x}_{x=0}^{+\infty} - \xi^2\FT_\cosine[f](\xi).
		\end{align}
	\end{subequations}
	\item 积分
	\begin{subequations}
		\label{eq:mmp-half-line-primitives}
		\begin{align}
			\FT_\sine\biggb{\int_0^xf(\eta)\d\eta}(\xi) &= \frac1\xi\FT_\cosine[f](\xi);\\
			\FT_\cosine\biggb{\int_0^xf(\eta)\d\eta}(\xi) &= -\frac1\xi\FT_\sine[f](\xi).
		\end{align}
	\end{subequations}
	其中$\textstyle \intRplus f(x)\d x = 0$。
	\item 象函数微分
	\begin{subequations}
		\label{eq:mmp-half-line-frequency-derivatives}
		\begin{align}
			\hat f_\sine' &= \FT_\cosine[xf]&\hat f_\sine'' &= -\FT_\sine[x^2f]\\
			\hat f_\cosine' &= -\FT_\sine[xf]&\hat f_\cosine'' &= -\FT_\cosine[x^2f].
		\end{align}
	\end{subequations}
	\item 半轴伸缩（$k > 0$）\begin{subequations}
		\label{eq:mmp-half-line-dilation}
		\begin{align}
				\FT_\sine[f(kx)](\xi) = \frac1k\hat f_\sine\Bigp{\frac\xi{k}},\\
				\FT_\cosine[f(kx)](\xi) = \frac1k\hat f_\cosine\Bigp{\frac\xi{k}}.
			\end{align}
\end{subequations}
	\item 在半轴上连续施行两次同种变换，由反演公式得（$\xi > 0$）\begin{subequations}
		\label{eq:mmp-half-line-inversion}
		\begin{align}
				\FT_\sine[\hat f_\sine(x)](\xi) &= \frac\pi{2}f(\xi),\\
				\FT_\cosine[\hat f_\cosine(x)](\xi) &= \frac\pi{2}f(\xi).
			\end{align}
\end{subequations}
\end{enumerate}
\paragraph{Fourier变换函数表}
下表的衰减指数和二次相位公式取$a > 0$。广义函数各式中，$1/x$取Cauchy主值，$1/x^n$取由其逐次求导定义的有限部，$1/(\i\xi)$也取主值。
\begin{enumerate}
	\item $\FT[\delta(x)] = 1,$\qquad$\FT[\delta^{(n)}(x)] = (\i\xi)^n;$
	\item $\FT[1] = 2\pi\delta(\xi),$\qquad$\FT[x^n] = 2\pi\i^n\delta^{(n)}(\xi);$
	\item $\FT[\e^{\i ax}] = 2\pi\delta(\xi - a);$
	\item $\FT[\cos ax] = \pi\Bigp{\delta(\xi - a) + \delta(\xi + a)};$
	\item $\FT[\sin ax] = -\i\pi\Bigp{\delta(\xi - a) - \delta(\xi + a)};$
	\item 对二次相位的余弦与正弦，记$u = \xi^2/(4a)$，则
	\begin{subequations}
		\label{eq:mmp-quadratic-phase-transforms}
		\begin{align}
			\FT[\cos ax^2] &= \sqrt{\frac{\pi}{2a}}(\cos u + \sin u),\\
			\FT[\sin ax^2] &= \sqrt{\frac{\pi}{2a}}(\cos u - \sin u).
		\end{align}
	\end{subequations}
	\item $\FT[\e^{-a\abs x}] = 2a/(a^2 + \xi^2);$
	\item $\FT[\e^{-ax^2}] = \sqrt{\pi/a}\e^{-\xi^2/(4a)};$
	\item $\FT[\e^{-\i ax^2}] = \sqrt{\pi/a}\e^{\i(\xi^2/(4a)-\pi/4)};$
	\item $\FT[1/x] = -\i\pi\sgn\xi;$
	\item 对有限部$1/x^n$，一般式与$n = 2$的特例为
	\begin{subequations}
		\label{eq:mmp-finite-part-fourier}
		\begin{align}
			\FT[\frac1{x^n}] &= -\i\pi\frac{(-\i\xi)^{n-1}}{(n - 1)!}\sgn\xi,\\
			\FT[\frac1{x^2}] &= -\pi\xi\sgn\xi.
		\end{align}
	\end{subequations}
	\item $\FT[\abs x^\alpha] = -2\Gamma(\alpha + 1)\sin(\pi\alpha/2)\abs\xi^{-1-\alpha},\quad -1 < \alpha < 0;$
	\item $\FT[\Heaviside(x)] = 1/(\i\xi) + \pi\delta(\xi);$
	\item $\FT[\Lambda](\xi) = \sinc^2(\xi/2)$；
	\item $\FT[\sinc](\xi) = \pi\Pi(\xi/2)$。
\end{enumerate}
这里 $\sinc x = \sin(x)/x$，在 $x = 0$ 连续取值为 $1$。单位宽矩形窗 $\Pi(x)$ 在
$\abs x < 1/2$ 时取 $1$，在 $\abs x > 1/2$ 时取 $0$，两个端点取 $1/2$。
两个矩形窗的卷积等于它们平移后的重叠长度，给出三角窗 $\Lambda(x) = (\Pi \ast \Pi)(x) = \max(1 - \abs x,0)$。
直接积分得
\[
	\FT[\Pi](\xi) = \int_{-1/2}^{1/2}\e^{-\i\xi x}\d x = \sinc(\xi/2).
\]
卷积定理给出表中的平方。
把矩形窗加宽一倍得 $\FT[\Pi(x/2)](\xi) = 2\sinc\xi$，再用 $\FT[\hat f](\xi) = 2\pi f(-\xi)$ 即得最后一式；
按对称截断取极限时，$\xi = \pm1$ 处的变换值为 $\pi/2$。

Heaviside阶梯函数
\begin{equation}
	\label{eq:mmp-heaviside}
	\Heaviside(x) := \begin{cases}
	1,&x > 0\\
	0,&x \leqslant 0
	\end{cases}
\end{equation}
\subsection{Laplace变换}
\paragraph{Laplace变换的性质}注：$x < 0$函数值取0。
\begin{enumerate}
	\item 线性
	\item 微分
	\begin{subequations}
		\label{eq:mmp-laplace-derivatives}
		\begin{align}
			\LT[f'] &= \xi\bar f - f(0);\\
			\LT[f''] &= \xi^2\bar f - \xi f(0) - f'(0);\\
			\LT[f^{(n)}] &= \xi^n\bar f - \Bigp{\xi^{n-1}f(0) + \xi^{n-2}f'(0) + \cdots + f^{(n-1)}(0)}.
		\end{align}
	\end{subequations}
	\item 积分 
	\begin{equation}
		\label{eq:mmp-laplace-primitive}
		\LT\biggb{\int_0^xf(\eta)\d\eta} = \frac1\xi\bar f.
	\end{equation}
	\item 像函数微分
	\begin{equation}
		\label{eq:mmp-laplace-frequency-derivative}
		\bar f' = \LT[-xf].
	\end{equation}
	\item 像函数积分
	\begin{equation}
		\label{eq:mmp-laplace-frequency-primitive}
		\int_\xi^{+\infty}\bar f(\xi)\d\xi = \LT\Bigb{\frac fx}.
	\end{equation}
	积分路径$\Re\xi > \sigma_0$。
	\item 复频移与时滞
	\begin{subequations}
		\label{eq:mmp-laplace-shift}
		\begin{align}
			\LT[f(x)\e^{\alpha x}] &= \bar f(\xi - \alpha),\quad \alpha \in \CC\\
			\LT[f(x - a)] &= \bar f\e^{-a\xi},\qquad\; a \geqslant 0 % \Heaviside(x-a)
		\end{align}
	\end{subequations}
	\item 相似
	\begin{equation}
		\label{eq:mmp-laplace-dilation}
		\LT[f(kx)] = \frac1k\bar f\Bigp{\frac\xi{k}},\quad k > 0
	\end{equation}
	\item 卷积
	\begin{equation}
		\label{eq:mmp-laplace-convolution}
		\LT[f \ast g] = \LT f\LT g
		\end{equation}
	\item 像函数卷积
	\begin{equation}
		\label{eq:mmp-laplace-product}
		\LT[fg](\xi) = \frac1{2\pi\i}\int_{\sigma-\i\infty}^{\sigma+\i\infty}\bar f(\zeta)\bar g(\xi - \zeta)\d\zeta,
		\end{equation}
	其中竖直积分线应同时位于两变换的收敛域，即$\sigma_f < \sigma < \Re\xi - \sigma_g$，并假定反演积分收敛。
\end{enumerate}
\paragraph{Laplace变换函数表}
\begin{enumerate}
	\item $\LT[\delta(t)] = 1,\quad \LT[\delta^{(n)}(t)] = s^n;$
	\item $\LT[\sgn t] = \LT[\Heaviside(t)] = 1/s,\quad \LT[t\Heaviside(t)] = 1/s^2;$
	\item $\LT\bigb{1/\sqrt{\pi t}} = 1/\sqrt{s};$%,\quad\LT\Bigb{2\sqrt{\frac t\pi}}=\frac1{s\sqrt s};$
	\item $\LT[\e^{at}] = 1/(s - a);$
	\item $\LT[\sin at] = a/(s^2 + a^2),\quad \LT[\sinh at] = a/(s^2 - a^2);$
	\item $\LT[\cos at] = s/(s^2 + a^2),\quad \LT[\cosh at] = s/(s^2 - a^2);$
\end{enumerate}

